\chapter{Introdução}\label{introduuxe7uxe3o}

A expansão das Infraestruturas de Medição Avançada (\emph{Advanced Metering Infrastructure} - AMI) altera de forma substancial a quantidade e a natureza da informação disponível às distribuidoras de energia. Em vez de depender apenas de leituras mensais de faturamento, passa a ser possível observar curvas de energia em intervalos de minutos, grandezas elétricas instantâneas, registradores, alarmes do equipamento e informações relacionadas à disponibilidade da própria comunicação. Esse aumento de resolução cria uma oportunidade relevante para reconhecer mudanças de comportamento, localizar eventos anômalos e fornecer novas evidências para investigação de perdas não técnicas.

Entretanto, maior disponibilidade de dados não implica automaticamente maior capacidade de identificar irregularidades. A detecção de perdas não técnicas é um problema particularmente difícil porque diferentes fenômenos podem produzir assinaturas semelhantes. Uma queda de energia registrada pode decorrer de redução legítima de carga, mudança de rotina, substituição do medidor, indisponibilidade de telemetria, geração local, defeito elétrico ou manipulação da medição. De forma análoga, uma UC com comportamento estatisticamente raro não é necessariamente uma UC irregular. A inspeção de campo permanece, portanto, como mecanismo de confirmação operacional.

A literatura de detecção de perdas não técnicas tem explorado técnicas supervisionadas e não supervisionadas, bem como métodos de ordenação de consumidores para inspeção. Sau et al.~tratam explicitamente a identificação e o \emph{ranking} de consumidores suspeitos como forma de aumentar a utilidade das inspeções \autocite{sau2023}. Trabalhos recentes também aproximam a detecção de furto da decisão operacional, incorporando valor econômico e restrições de inspeção \autocite{han2026}, enquanto abordagens explicáveis buscam reduzir a distância entre o escore produzido pelo modelo e a interpretação realizada pelo especialista \autocite{pavlidis2026}.

Neste trabalho, localização designa a associação do desvio a uma UC, a um período e, quando disponível, a um alimentador. O objetivo central deste trabalho é a \textbf{detecção, localização e caracterização de padrões anômalos em redes de distribuição}, com interesse especial em padrões compatíveis com perdas não técnicas e furto de energia. A priorização de inspeções constitui aplicação operacional da detecção e da caracterização das evidências. Para apoiar essa triagem, o sistema desenvolvido entrega ao analista a identificação da UC, o período em que o desvio foi observado, a condição da comunicação e os indicadores elétricos ou de alarme que motivaram a seleção, permitindo confrontar o volume listado com a capacidade hipotética de vistoria.

A formulação inicial do trabalho partiu de uma abordagem global de detecção de anomalias: janelas históricas eram transformadas em atributos estatísticos, e diferentes detectores não supervisionados eram combinados para produzir um escore de anomalia. Essa ordenação global serviu de referência comparativa (\emph{baseline}) nos experimentos deste trabalho. A inspeção das séries AMI mostrou, contudo, que lacunas de telemetria, trocas de medidor, mudanças de patamar e diferenças entre alimentadores alteravam simultaneamente o mesmo vetor de atributos, fazendo com que falhas de comunicação ou alterações cadastrais recebessem pontuação semelhante à de desvios elétricos.

A arquitetura atual surge dessa observação. Em vez de substituir o objetivo original, ela decompõe o problema em visões complementares: qualidade e comunicação, mudança temporal, comportamento diário, grandezas elétricas, alarmes físicos, regras determinísticas e contexto topológico. Essa decomposição permite testar separadamente a qualidade da telemetria, a raridade elétrica e a contextualização temporal ou topológica. A comparação retrospectiva examina como cada componente modifica as hipóteses e o ranking, sem atribuir a essas mudanças, por si, redução comprovada de falsos positivos.

Outro ponto central é a natureza temporal da própria identidade do dado. A telemetria é indexada pelo medidor físico, enquanto a unidade de interesse comercial e operacional é a Unidade Consumidora (UC). Uma UC pode receber diferentes medidores ao longo do tempo; um mesmo medidor pode ser retirado, substituído ou reutilizado. Sem reconstruir essa relação por vigência temporal, uma troca de equipamento pode produzir uma falsa quebra de consumo e ser confundida com uma anomalia real. Esse problema de engenharia de dados passa, portanto, a fazer parte da metodologia de detecção e não apenas da etapa de ETL.

A escala do problema também exige uma solução de software compatível com a realidade operacional. A aplicação desenvolvida não se restringe aos modelos de Machine Learning. Ela inclui extração incremental e idempotente, Data Lake particionado, processamento colunar com Polars, escrita por \emph{sinking}, auditoria de continuidade temporal, controle de integridade, geração de datasets, treinamento, versionamento de modelos e interfaces para exploração dos resultados. A TUI organiza os fluxos de engenharia e governança, enquanto o protótipo funcional do Cockpit Web permite navegar pelas evidências produzidas para cada UC. O vínculo com ordens de serviço e inspeções de campo não integra a implementação avaliada.

Essa evolução também oferece uma base para a transformação futura do processo operacional. Nas execuções retrospectivas, o sistema produz listas dinâmicas de evidências; em uma implantação, essas listas poderiam ser direcionadas à revisão humana e, conforme decisão da distribuidora, a filas de comunicação, manutenção ou eventual investigação comercial. A hipótese diagnóstica não determina automaticamente o encaminhamento. O número de casos produzidos pode então ser confrontado, em cenário de planejamento, com uma capacidade hipotética de inspeção. Essa simulação examina a organização possível da carteira; não representa execução de ordens de serviço ou medição de rendimento em campo.

Neste documento, a reconstrução temporal das séries AMI e o comportamento das listas de seleção correspondem a resultados computacionais avaliados sobre dados reais, ao passo que a governança de filas e o piloto de campo delimitam uma aplicação futura.



